\documentclass[aps,notitlepage]{revtex4-2}
\usepackage[utf8]{inputenc}
\usepackage[margin=1in]{geometry}
\usepackage{microtype}
\usepackage{amsmath, amsthm, amssymb, amsfonts}
\usepackage{physics}
\usepackage{xcolor}
\usepackage{graphicx}
\usepackage{natbib}
\usepackage{hyperref}
\usepackage{todonotes}
\usepackage{subcaption}
\usepackage[section]{placeins}

    % \setlength{\parindent}{0em}

    \newcommand{\cbqty}[1]{\left\{#1\right\}}
    \newcommand{\nts}[1]{\textcolor{blue}{#1}}
    \newcommand{\fixme}[1]{\textcolor{red}{#1}}

    \newcommand{\sigx}{\sigma^x}
    \newcommand{\sigy}{\sigma^y}
    \newcommand{\sigz}{\sigma^z}

\begin{document}

\title{QUBO encoding effects on quantum annealing efficiency}
\date{\today}
\maketitle

\section{The JobShop problem}

Referring to scheduling theory~\cite{pinedo2012scheduling} we consider the \emph{JobShop} $(J_m)$ problem with release and due constraints ($r_i d_i$), where the objective is total weighted tardiness $\sum_j w_j T_j$, namely:
\begin{equation}
    J_m | r_i d_i | \sum_j w_j T_j.
\end{equation}
In the \emph{JobShop} each job has its own schedule in terms of the sequence of machines it has to process on, and each machine can process one job at a time (i.e. sometimes jobs have to wait in the buffer before the machine is free). We consider the release constraint as the earliest possible time the job can be submitted to the first machine, while the die constraint is the time the job has to be performed on all machines. Beyond the due constraint, each job has to be performed as soon as possible on the sequence of machines, to minimize the objective of weighted tardiness (the time the job is processed on the last machine).
Such a scheduling problem has many practical applications and is generally NP-Hard~\cite{baptiste2001constraint}.

To keep up with standard linear encoding of such problem~\cite{pinedo2012scheduling} we use the following definitions: 
\begin{itemize}
\item Let $S_j$ be the schedule of job $j$, i.e. the series of machines it has to be processed on,
\item Let $p_{j,m}$ be the processing time of job $j$ on machine $m$. 
\end{itemize}
These two, together with release times $r_j$, die times $d_j$ and weights $w_j$ fully define our\emph{JobShop} problem. For the sake of encoding, we need also the following definitions.
\begin{itemize}
\item Let $\mathcal{S}_{j,j'}$ a set of machines used both by $j$ and $j'$. 
\item Let $\sigma(S_j, m)$ be the preceding machine  to $m$ in the schedule $S_j$. 
\item Let $m_{j, \text{start}}$ and $m_{j, \text{end}}$ be the first and last machine of job $j$. 
\item Let $S^+_j(m)$ be the schedule of job $j$, up to machine $m$ (including $m$) and $S^-_j(m)$ the reminding part of the schedule such that joining of $S^+_j(m)$ and $S^-_j(m)$ produces $S_j$.
\end{itemize}

\subsection{Integer Linear Programming encoding}

The state-of-art Integer Linear Programming (ILP) encoding with precedence variables requires the following decision variables:
\begin{itemize}
    \item $t_{j,m} \in \mathbb{Z}^+$ - time job $j$ is finished on machine $m$,
    \item $y_{j, j', m} \in \{0,1\}$ - precedence variable equal to one if $j$ is performed before $j'$ on $m$.
\end{itemize}
(Obviously, we can use $t \in \mathbb{R}^+$ yielding mixed integer linear programming approach, this however would impose complications in the QUBO version of the problem).

From \emph{Job Shop} definition, each job has to be processed on a sequence of machines and one machine can handle only one job at a time, yielding the following constraints:
\begin{equation}
    \begin{split}
      \forall_{j} \forall_{m \in S_j^- (m_{j, \text{start}})} \ \ \  t_{j,\sigma(S_j, m)} + p_{j,m} &\leq t_{j,m}  \\
     \forall_{j \neq j'} \forall_{m \in \mathcal{S}_{j,j'}} t_{j', m} + p_{j,m} &\leq M y_{j,j', m} + t_{j,m} 
    \end{split}
    \label{eq::ILP_machines}
\end{equation}
In the second inequality $M$ is the high enough number (in co-called big-M encoding) to always hold if $y_{j,j',m} = 1$. (Minimal value of such $M$ can be obviously computed by using minimal value of $t_{j,m}$ and maximal value of $t_{j', m}$) discussed later.

The release and due constraint yields:
\begin{equation} 
\begin{split}
    \forall_{j} \  \ \ r_j + p_{j, m_{j, \text{start}}} &\leq t_{j, m_{j, \text{start}}} \\
    \forall_{j} \ \ \  t_{j, m_{j, \text{end}}} &\leq d_j
\end{split}\label{eq::ILP_constr_r_d}
\end{equation}
From Eq.~\eqref{eq::ILP_constr_r_d} we can conclude the minimal and maximal values of each variable:
\begin{equation} 
\begin{split}
    &\forall_{j}  \forall_{m \in S_j}\ \ \ \ t_{\text{min}}(j,m)  = r_j + \sum_{m' \in S^+_j(m)} p_{j, m'} \leq t_{j, m} \\
    &\forall_{j} \forall_{m \in S_j}\ \ \ \ t_{\text{max}}(j,m) =  d_j - \sum_{m' \in S^-_j(m)} p_{j, m'} \geq t_{j, m} 
\end{split}\label{eq::ILP_constr_limits}
\end{equation}

The objective of maximal weighted tardiness can be represented as follows:
\begin{equation}
    \text{objective} =  \sum_j w'_j t_{j, m_{j, \text{end}}} - \text{offset} 
    \label{eq::ilp_obj}
\end{equation}
where 
\begin{equation}
\begin{split}
    w'_j &=  \frac{w_j}{t_{\max}(j, m_{j, \text{end}}) - t_{\min}(j, m_{j, \text{end}}) }  \\
    \text{offset} &= \sum_j w'_j t_{\min}(j, m_{j, \text{end}}) .
\end{split}
\end{equation}
The linear scaling does not matter for linear constrained problems but is used bet better comparison with following QUBO encoding. In particular, input from each $ t_{j, m_{j, \text{end}}}$ does not exceed $w_j$ which is helpful in the selection of constraints penalties in QUBO encoding more straightforward.

\subsection{QUBO encoding}

Such representation is easily convertible to QUBO.
Following \cite{venturelli2016job}, we derive QUBO directly from \emph{JobShop} problem by using following decision binary variables 
\begin{equation}
    x_{j,m,t} = \begin{cases} 1 \text{ if job } j \text{ is completed on machine } m \text{ at time } t \\
    0 \text{ otherwise},
    \end{cases}
\end{equation}
To enforce release and due constraint we limit each $t$ index by its minimal and maximal value in Eq.~\eqref{eq::ILP_constr_limits}
an ensure that each job is completed once and only once:
\begin{equation}\label{eq::sum_to_one_constraint}
  \forall_j \forall_{m \in S_j} \sum_{t_{\text{min}}(j,m) \leq t \leq t_{\text{max}}(j,m)} x_{j,m,t} = 1,
\end{equation}
all available binary variables can be then concatenated into the vector $\vec{x}$.

Constraint in Eq.~\eqref{eq::sum_to_one_constraint} can be transformed to the unconstrained version, by the penalty method~\cite{gusmeroli2022expedis}, leading to the minimization of the form of Eq.~\eqref{eq::sum_to_one}.

Eq.~\eqref{eq::ILP_machines} is encoded in QUBO as follows:
\begin{equation}
    \begin{split}
      \forall_{j} \forall_{m \in S_j^- (m_{j, \text{start}}) }  \forall_{t_{\text{min}}(j,\sigma(S_j, m)) \leq t' \leq t_{\text{max}}(j,\sigma(S_j, m))} \forall_{t_{\text{min}}(j,m) \leq t < t' + p_{j,m} }\ \ \  x_{\sigma(S_j, m),j,t'} x_{m, j, t} &= 0 \\
     \forall_{j \neq j'} \forall_{m \in \mathcal{S}_{j,j'}} \forall_{t_{\text{min}}(j,m) \leq t \leq t_{\text{max}}(j,m)}  \forall^*_{t - p_{j,m} < t' < t + p_{j',m}} \ \ \ \  x_{j,m,t} x_{j', m, t'} &= 0
    \end{split}
    \label{eq::QUBO_machines}
\end{equation}
* - obviously we also require $t_{\text{min}}(j',m) \leq t' \leq t_{\text{max}}('j,m)$.
Constraints in Eq.~\eqref{eq::QUBO_machines} yields unconstrained problem as in Eq.~\eqref{eq::pair}.

The objective from Eq.~\eqref{eq::ilp_obj} is implemented as the following QUBO term:
\begin{equation}
    \text{objective} = \text{min}_{\vec{x}}  \sum_j \sum_{t_{\min}(j, m_{j, \text{end}}) \leq t \leq t_{\max}(j, m_{j, \text{end}})} w'_j t   x_{j, m_{j, \text{end}}, t} - \text{offset}
    \label{eq::ilp_obj}
\end{equation}
(Observe that dislike the ILP approach, generalization to non-linear objective in $t$ is straightforward).


\section{QUBO parameters}

When we convert a constrained optimization problem into QUBO form, typically we introduce penalty parameters to ensure that the solution of the unconstrained optimization problem actually satisfies all the original constraints.
\begin{align}
\label{eq::sum_to_one}
    p_{\rm sum}\sum_{m,j}\pqty{\sum_{t,t'} x_{m,j,t}x_{m,j,t'} - \sum_{t} x_{m,j,t}^2}
\end{align}
and (in symmarised version)
\begin{align}
\label{eq::pair}
    p_{\rm pair}\sum_{i<i'}\pqty{x_i x_{i'} + x_{i'} x_i}
\end{align}
where indices $i$, $i'$ are derived from corresponding $m$, $j$, $t$.
The specific choice of penalty parameters can have a significant and non-obvious impact on how well any given optimizer performs on the resulting problem as we have seen in the trains work, for example comparing the results obtained from the QUBOs generated with different penalty values leading to split and unsplit spectra.

If we consider optimization with quantum annealing specifically, we can take a thermodynamic perspective on the optimization process, e.g. by studying the computational and thermodynamic efficiency as in \cite{Smierzchalski2024}. 
Broadly, the idea to pursue is to apply these ideas to the problem of setting values for penalties when solving a constrained optimization problem using quantum annealing, and in understanding different optimization regimes. 

For example, take some particular QUBO from the trains work. 
We could ask:
\begin{itemize}
    \item How do the efficiencies $\eta_{\rm th}$, $\eta_{\rm comp}$, and $Q_{\rm GS}$ vary as $p_{\rm sum}$ and $p_{\rm pair}$ are varied? Is there any kind of ``transitional'' behavior in any of these quantities corresponding to split vs. unsplit spectra in the original QUBO? Does this change with the annealing time?
    \item It seems that on the real device, reverse-annealing is the most accessible way to approach these questions. Is there some way to get enough information from forward annealing?
    \item Do results with reverse-annealing inform results with forward annealing?
\end{itemize}

This could also expand to include different problem encodings, for example comparing precedence variables against the discrete approach to time-stepping used in the trains work.

\section{Main question}

When we vary the hyperparameters of our optimization problems (binary encoding, penalty parameters, annealing details, etc.), what happens to the efficiency of the optimization process? Can this tell us anything about how to choose these hyperparameters, or show us some interesting differences e.g. in the split vs. overlapping cases?


We can think about solving QUBOs on the classical device and asses the 
spectrum as well as quality of solution in this way.

\section{Numerical studies}

Numerical simulations seem like a reasonable tool for this, as they should be able to give more detailed results through fine parameter sweeps without needing to run a huge number of experiments on a real annealer.

Simulating a noisy quantum annealer can be done for small scale problems by treating it as an open quantum system and running master equation simulations. 
In particular, the master equations derived in \cite{Albash2012} for describing ``\emph{the time evolution of a system evolving adiabatically while coupled weakly to a thermal bath}'' seem to be a reasonable starting point. These constructions have been used to study quantum annealing in a variety of works, e.g. for studying the effects of pausing in annealing \cite{Chen2020}. 
That said, there are limitations to the technique, e.g. \cite{Bando22}, but it should hopefully suffice for problems roughly up to the size of what we were able to solve with QAOA on the IonQ device.

In \cite{Albash2015b} this type of master equation was used alongside two potential non-quantum models of a D-Wave machine to validate some earlier claims about witnessing multi-qubit entanglement during the annealing process. 
One important point is made in the conclusion of this work, the authors point point out that ``\emph{it becomes computationally prohibitive to use the ME to study systems larger than about 15 qubits}''. 
15 qubits is certainly not large, but it is also not trivially small so it should hopefully be possible to draw meaningful conclusions from such simulations.

\subsection{Approach}

Once we have a master equation simulation of the quantum annealing process, we can study fluctuation relations and compute the thermodynamic and computational efficiencies as has been done before experimentally \cite{Smierzchalski2024, Campisi2021, Buffoni2020, Gardas2018}. 
We can bypass several of the limitations of real experiments, for example in addition to reverse annealing (with or without a pause) we can study the standard forward annealing process with a known initial state.
Also, if the simulation is detailed enough then we can read out or at least estimate thermodynamic quantities such as the entropy production, work, and heat flux which should allow us to determine the efficiency without having to rely on thermodynamic uncertainty relations or other inequalities. 

Ultimately the goal of the simulation work is to have a code which takes a QUBO/Ising Hamiltonian $H$ along with a description of the annealing process (annealing schedule $A(t)$ \& initial state $\psi_0$) and computes the thermodynamic and computational efficiencies, $\eta_{\rm th}$ and $\eta_{\rm comp}$, as well as the (normalized) average solution energy $Q_{\rm GS}$.

\subsection{What to simulate}

For small QUBOs, hopefully the simulation process is quick enough that we can sweep various parameters and study how the efficiencies vary. 
These might be good things to run on the HPC resources we have here at UMBC.

There are a few plots which seem like clear targets, first with forward annealing:
\begin{enumerate}
    \item Sweep the penalty parameters $p_{\rm sum}$ and $p_{\rm pair}$. Do either of the efficiencies show any kind of transition when the QUBO spectrum is split vs. unsplit? For this, we would also want to diagonalize the QUBOs exactly and plot some kind of signed gap between feasible and unfeasible solutions.
    \item For a couple choices of QUBO parameters (e.g., split vs. unsplit vs. boundary), how do the efficiencies scale with the annealing time? Is there any difference in scaling in the different regimes?
    \item These experiments can be repeated for the two types of master equations constructed in \cite{Albash2012}, to see if the results are robust predictions or if they are highly dependent on the exact model used. 
\end{enumerate}
and also with reverse annealing:
\begin{enumerate}
    \item Repeat the same simulations as with forward annealing but with reverse annealing, possibly including a pause. Are the efficiency results related to the forward annealing case? If they are, this might mean that on the real device we can learn about the behavior of forward annealing from reverse annealing which seems easier to characterize.
    \item Repeat the same simulations (or build this into the previous) with reverse annealing only giving access to the change in energy, just like the real device. This allows checking how accurate the estimates for the efficiencies are when only partial information is available and thermodynamic uncertainty relations are used to fill in the rest. 
    \item Is there some easily-computed initial guess state which makes reverse annealing give a good solution to the QUBO? If so, how do the efficiencies compare to forward annealing?
    \begin{itemize}
        \item[-] E.g., maybe take the ``linear'' part of the QUBO only and start with a thermal state defined from there with some variable $\beta$.
    \end{itemize}
\end{enumerate}
and finally with the different translations of a scheduling problem into a QUBO:
\begin{enumerate}
    \item Assuming that formulation also has some penalty parameters, repeat the same analysis for the new formulation and compare. 
\end{enumerate}

\subsection{Additional possibilities}

Several works have looked at the ``quantum-ness'' of the D-Wave machine by checking if there is some classical annealing-like model which can reproduce features of the output, for example \cite{Albash2015b}. 
That work looked at whether the output of the D-Wave machine can be explained through a simulated quantum annealing (SQA) model or if it could be explained as the output of a particular classical rotor model (SSSA).
These are still physical models which have the vague structure of a quantum annealer (a collection of degrees of freedom subject to some programmable coupling), but the dynamics they run are something else.

Whatever analysis we do with master equation simulations of quantum annealing could be repeated for one or more of these classical models.
The goal would be to check which if any of the qualitative efficiency features of the classical model match with the expectations for a quantum annealer. 
E.g., maybe the fluctuation relations for the different underlying models translate into some difference in the scaling of the efficiency with the annealing time.
Alternately, if the qualitative behavior of the efficiency is captured by a classical model, we could use that to check how the efficiency should behave for larger problems.

This would be something to think about once the primary study has been done.

\section{Experiments}

Exactly what might be worth testing on the real device will depend on what the simulations show.
Some possibilities might include:
\begin{itemize}
    \item Do a coarse 1D sweep in the space of penalty parameters and check that the real device scales like the simulations. Do a coarse 1D sweep in the annealing time, also for verification.
    \item If our simulations show that there is some relationship between the behavior of forward and reverse annealing, do the above sweeps with both schedules and check if the device appears to exhibit the same relationship.
    \item Look at the scaling with problem size vs. qubit count, since this won't be tractable with simulations.
\end{itemize}
The simulations definitely will not have the fidelity required to match in detail with the device output, but hopefully the qualitative behavior can match.

\section{Gate-based extensions}

There is nothing in the definitions of the thermodynamic and computational efficiencies that inherently ties them to quantum annealing.
If we had a way to compute the underlying average thermodynamic quantities for say a quantum circuit which implements QAOA, we could just as easily compute the efficiencies for a gate-based approach to solving QUBOs.

It might be interesting then to see if the work Josey is doing on the energetic costs required to implement quantum gates could be used to estimate or bound the work input required to implement the QAOA ansatz circuit. 
If so, this can possibly be built into a procedure to determine the efficiencies for QAOA at some point in the parameter space.
At that point, all the same tests as with quantum annealing could be performed, as well as some new ones:
\begin{enumerate}
    \item What is the impact of additional layers in the ansatz on the efficiency?
    \item The efficiencies are measured for some particular set of parameters for the parameterized circuit $\vec{\beta}$ and $\vec{\gamma}$. How does the efficiency vary as these parameters are changed, with all the problem parameters fixed? Is the efficiency at the optimal point uniquely good (or bad)?
\end{enumerate}

% \bibliographystyle{plain}
\bibliography{references}

\end{document}
